\documentclass[11pt,a4paper]{report}
\usepackage[T1]{fontenc}
\usepackage[utf8]{inputenc}
\usepackage[margin=1in]{geometry}
\usepackage{amsmath,amssymb}
\usepackage{booktabs}
\usepackage{longtable}
\usepackage{xcolor}
\usepackage{fancyvrb}
\usepackage{titlesec}
\usepackage{enumitem}
\usepackage[hidelinks,bookmarksnumbered]{hyperref}
\usepackage{parskip}

\definecolor{codebg}{gray}{0.96}
\definecolor{rule}{gray}{0.75}
\DefineVerbatimEnvironment{Rblock}{Verbatim}%
  {fontsize=\small, frame=leftline, framerule=1.2pt,
   rulecolor=\color{rule}, framesep=10pt, xleftmargin=12pt,
   samepage=false}

\titleformat{\chapter}[hang]{\LARGE\bfseries}{\thechapter.}{0.6em}{}
\titlespacing*{\chapter}{0pt}{0pt}{18pt}
\setcounter{secnumdepth}{2}
\setcounter{tocdepth}{1}

\title{\vspace{-1.5cm}\textbf{\Huge tpridge}\\[6pt]
{\large Two-Parameter Ridge Estimation of\\HAR Realized-Volatility Models}\\[14pt]
{\normalsize User Manual, version 0.2.0}}
\author{Talha Omer \and Muhammad Ibrar Khan \and Kristofer M\aa nsson \and
P\"ar Sj\"olander}
\date{}

\begin{document}
\maketitle
\thispagestyle{empty}
\vfill
\noindent\small
This manual is written to be followed from start to finish. Every block
of code runs as printed, and the chapters build on one another: after
Chapter 3 you will have estimated the model on simulated data, after
Chapter 4 on a classical ill-conditioned regression, after Chapter 6 on
real cryptocurrency data downloaded inside R, and after Chapter 8 you
will have reproduced the tables of the accompanying paper.

\medskip
\noindent
Chapter 4 is the one to read if you came for the estimator rather than
for volatility. The two-parameter ridge is a general remedy for an ill
conditioned design, and \texttt{tpr\_lm()} applies it to any
regression.

\medskip
\noindent
The reference manual, generated from the sources, documents every
argument of every function. Build it with
\texttt{Rscript build\_manual.R}. This document is the companion to it:
it explains what to do and why, rather than listing arguments.
\newpage

\tableofcontents
\newpage

% =====================================================================
\chapter{Installation}

\section{From GitHub}

\begin{Rblock}
install.packages("remotes")
remotes::install_github("talhaomer/tpridge", build_vignettes = TRUE)
library(tpridge)
\end{Rblock}

\section{From a local copy}

If you have the source directory rather than the repository:

\begin{Rblock}
install.packages("devtools")
devtools::install("path/to/tpridge", build_vignettes = TRUE)
\end{Rblock}

\section{What it depends on}

The package needs \texttt{data.table} and nothing else to run. Four
functions ask for a further package only when you call them, and tell
you so if it is missing:

\begin{center}
\begin{tabular}{ll}
\toprule
Function & Needs \\
\midrule
\texttt{tpr\_mcs()} & \texttt{MCS} \\
Parallel Monte Carlo & \texttt{future}, \texttt{future.apply} \\
Vignette building & \texttt{knitr}, \texttt{rmarkdown} \\
Tests & \texttt{testthat} \\
\bottomrule
\end{tabular}
\end{center}

To install everything at once:

\begin{Rblock}
install.packages(c("data.table", "MCS", "future", "future.apply",
                   "knitr", "rmarkdown", "testthat"))
\end{Rblock}

\section{Checking the installation}

\begin{Rblock}
library(tpridge)

set.seed(1)
X <- matrix(rnorm(400), ncol = 4)
y <- as.numeric(X %*% c(1, 0.6, 0.3, 0.1)) + rnorm(100)
fit <- tpr_fit(X, y)
fit
\end{Rblock}

You should see a block reporting the penalty, \texttt{q}, the effective
degrees of freedom, the three coefficients of determination, and a line
reading \emph{Proposition 2 residual} followed by a number of order
$10^{-16}$. That last line is the package verifying its own central
identity; if it is not tiny, something is wrong and you should not
proceed.

% =====================================================================
\chapter{The idea in one page}

Ridge regression adds a penalty $\lambda$ to the normal equations,
$$\hat\beta_{I}(\lambda) = (X'X + \lambda I)^{-1} X'y .$$
This reduces variance and is why ridge helps when regressors are
collinear. It also fits the estimation sample worse than least squares,
necessarily so, because least squares maximises the in-sample fit by
construction.

The two-parameter estimator rescales the ridge vector by a scalar,
$$\hat\beta_{II}(\lambda) = q(\lambda)\,\hat\beta_{I}(\lambda),
\qquad
q = \frac{\hat\beta_I' r}{\hat\beta_I' G \hat\beta_I},
\qquad G = X'X, \; r = X'y,$$
choosing $q$ to maximise the fit along the ray $\{q\hat\beta_I\}$. Four
results follow, and the package verifies each of them numerically.

\paragraph{The scalar always exceeds one.}
$q = 1 + \lambda\|\hat\beta_I\|^2/\|X\hat\beta_I\|^2 > 1$ for every
$\lambda>0$. It is derived, not assumed.

\paragraph{The fit gain has a closed form.}
$$R^2(\hat\beta_{II}) - R^2(\hat\beta_{I})
= \frac{(q-1)^2\,\|X\hat\beta_I\|^2}{y'y}.$$
This is the central result. It says exactly how much fit the second
parameter recovers, in terms of quantities already computed.

\paragraph{Nothing beats least squares in sample.}
$R^2(\hat\beta_{OLS}) \ge R^2(\hat\beta_{II}) \ge R^2(\hat\beta_{I})$
always. The second parameter closes part of a gap; it never crosses it.
Any claim of an in-sample improvement over least squares would be false.

\paragraph{The two estimators fail differently.}
As $\lambda\to\infty$ the one-parameter estimator collapses to the null
predictor, while the two-parameter estimator converges to a finite
non-zero limit. Under a badly chosen penalty this difference is large,
and since the right penalty is not observable in advance, it is the
practical argument for the method.

\section{A warning about $\lambda$}

The penalty is not scale free. The eigenvalues of $X'X$ grow with the
sample size, so the same $\lambda$ delivers different amounts of
shrinkage in different designs, and reporting ``$\lambda = 2000$'' tells
the reader nothing. The package therefore works with the effective
degrees of freedom,
$$\mathrm{df}(\lambda) = \sum_{j=1}^{p} \frac{e_j}{e_j + \lambda},$$
which falls from $p$ at $\lambda = 0$ towards zero. Saying ``the model
retains 1.5 of 3 effective parameters'' is interpretable; saying
$\lambda = 2000$ is not.

% =====================================================================
\chapter{First steps: the estimators alone}

The three core functions take a design matrix and a response. They know
nothing about volatility and can be used on any regression problem.

\section{Choosing a penalty}

\begin{Rblock}
set.seed(1)
X <- matrix(rnorm(400), ncol = 4)
y <- as.numeric(X %*% c(1, 0.6, 0.3, 0.1)) + rnorm(100)

tpr_lambda(X, y)                                # data-driven, HKB rule
tpr_lambda(X, y, rule = "df", target_df = 2)    # penalty for a target df
\end{Rblock}

The first is the rule of Hoerl, Kennard and Baldwin (1975),
$\hat\lambda = p\hat\sigma^2 / \hat\beta'_{OLS}\hat\beta_{OLS}$. The
second solves $\mathrm{df}(\lambda) = 2$ for $\lambda$, so you state the
shrinkage you want and the function finds the penalty that delivers it.

\section{The scalar on its own}

\begin{Rblock}
q <- tpr_q(X, y, lambda = 20)
q                       # always greater than one
attr(q, "eff_df")       # the shrinkage this corresponds to
attr(q, "rho")          # the Rayleigh quotient that determines q
\end{Rblock}

\section{Fitting all three estimators}

\begin{Rblock}
fit <- tpr_fit(X, y, lambda = 20)
fit

coef(fit, "OLS")
coef(fit, "Ridge-I")
coef(fit, "Ridge-II")

fit$r2_in          # in-sample fit of all three
fit$recovery       # share of the sacrificed fit that is recovered
fit$condition      # condition number of the design
\end{Rblock}

\section{Seeing the trade-off}

\begin{Rblock}
plot(fit, which = "path")   # R2 against effective degrees of freedom
plot(fit, which = "coef")   # the three coefficient vectors
plot(fit, which = "gap")    # the fit gain against the penalty
\end{Rblock}

The first plot is the one to look at. Least squares is flat, because it
does not depend on the penalty. As shrinkage increases, the
one-parameter estimator falls away steeply while the two-parameter
estimator stays close to least squares. The vertical dotted line marks
the penalty actually selected.

\section{Verifying the theory yourself}

\begin{Rblock}
for (l in c(1, 20, 500, 5000, 1e6)) {
  f <- tpr_fit(X, y, lambda = l)
  cat(sprintf("lambda %8.0f  q %7.4f  residual %.2e  ordering %s\n",
      l, f$q, abs(f$gap_empirical - f$gap_theoretical),
      f$r2_in[["OLS"]] >= f$r2_in[["Ridge-II"]] &&
      f$r2_in[["Ridge-II"]] >= f$r2_in[["Ridge-I"]]))
}
\end{Rblock}

Every residual should be of order $10^{-16}$, every $q$ above one, and
every ordering \texttt{TRUE}.

% =====================================================================
\chapter{Any regression, not only HAR}

Nothing about the two-parameter ridge estimator is specific to
volatility. It is a remedy for an ill-conditioned design, and the design
can come from anywhere. \texttt{tpr\_lm()} gives it the interface you
would expect, so it can be used the way \texttt{lm()} is used.

\section{The formula interface}

\begin{Rblock}
data(longley)
m <- tpr_lm(Employed ~ GNP + Unemployed + Armed.Forces +
              Population + Year, data = longley)
m
\end{Rblock}

Longley (1967) assembled sixteen annual observations on United States
employment against six macroeconomic series, as a test of the numerical
accuracy of regression software. The regressors are so nearly collinear
that early programs disagreed with one another on the answer. It is the
canonical ill-conditioned design and the one on which ridge regression
is usually demonstrated.

\section{Diagnostics}

\begin{Rblock}
summary(m)
\end{Rblock}

The summary prints the three coefficient vectors side by side, then the
diagnostics that bear on whether shrinkage is called for:

\begin{Rblock}
Collinearity diagnostics:
  condition number of X'X      6507.3
  smallest eigenvalue          0.008918
  variance inflation factors
    GNP                   1034.22  *
    Unemployed              23.26  *
    Armed.Forces             3.15
    Population             225.71  *
    Year                   732.64  *

Shrinkage:
  lambda 0.005037,  q 1.001307,  effective df 4.573 of 5
\end{Rblock}

A condition number in the thousands and variance inflation factors above
a thousand say the design is severely ill conditioned. Gross national
product, population and the year index are all measuring the same
expanding economy.

\section{What this example teaches}

The recovery figure in that summary is $0.4$ per cent, which looks like
a disappointment until one reads the line above it. The data-driven
penalty is $\lambda = 0.005$, giving $q = 1.0013$, and the fit gain is
proportional to $(q-1)^2$. The correction has almost nothing to work
with.

This is worth dwelling on, because the natural conjecture is the
opposite: that severe collinearity should make the second parameter
matter more. Proposition 1 bounds $q$ above by
$1 + \lambda/\lambda_{\min}(G)$, so ill-conditioning raises the
\emph{ceiling} on the attainable gain. Whether the gain is attained
depends on the realised $q$, which depends on $\lambda/\rho$, and a
data-driven penalty on a real design generally leaves that ratio small.

Fix the shrinkage instead of letting the data choose it, and the
mechanism becomes visible:

\begin{Rblock}
for (df in c(5, 4, 3, 2, 1)) {
  f <- tpr_lm(Employed ~ GNP + Unemployed + Armed.Forces +
                Population + Year, data = longley,
              rule = "df", target_df = df)
  cat(sprintf("df %4.1f   q %7.4f   recovery %5.1f%%\n",
              df, f$q, 100 * f$recovery))
}
\end{Rblock}

As the penalty tightens, $q$ grows and so does the share of the
sacrificed fit that the second parameter restores.

\section{Prediction and residuals}

\begin{Rblock}
predict(m, longley[1:3, ], estimator = "Ridge-II")
predict(m, longley[1:3, ], estimator = "OLS")
head(residuals(m, "Ridge-II"))
\end{Rblock}

\section{A word on inference}

\texttt{summary()} reports no standard errors, and this is deliberate
rather than an omission. The sampling distribution of a shrinkage
estimator is not that of least squares, so the usual standard errors,
$t$ statistics and confidence intervals do not carry over. Printing them
would invite a reader to interpret them as if they did.

If inference is needed, a bootstrap that resamples the entire estimation
procedure, penalty selection included, is the defensible route. For
least squares coefficients on the same data, \texttt{lm()} together with
a heteroskedasticity and autocorrelation consistent covariance estimator
remains appropriate.

\section{When is shrinkage worth it}

The diagnostics answer the first half of the question. A condition
number in the hundreds or thousands, or variance inflation factors above
ten, indicate a design on which least squares coefficients are unstable.

Whether the \emph{second} parameter adds anything is a separate question,
and the summary answers that too, through $q$. Close to one and the
correction is negligible; materially above one and it recovers most of
the fit that one-parameter ridge gives up.

% =====================================================================
\chapter{Downloading data}

\section{Five-minute prices from Binance}

\begin{Rblock}
k <- binance_klines("BTCUSDT", from = "2024-01", to = "2024-06")
head(k)
\end{Rblock}

No API key is needed: the monthly archives are public. Files are cached,
so a month already downloaded is not fetched again. To keep them between
sessions, give a directory:

\begin{Rblock}
k <- binance_klines("BTCUSDT", from = "2018-05", to = "2025-12",
                    dir = "~/binance_cache")
\end{Rblock}

The five pairs used in the paper:

\begin{Rblock}
binance_symbols()
#> "BTCUSDT" "ETHUSDT" "LTCUSDT" "XRPUSDT" "BNBUSDT"
\end{Rblock}

\section{Two details that matter}

Binance changed the \texttt{open\_time} field from milliseconds to
microseconds partway through the history. If you write your own reader
and do not handle this, every recent date will be wrong by a factor of a
thousand. \texttt{binance\_klines()} detects and normalises it. Some
archives also carry a header row while others do not, which is likewise
detected.

\section{Realized measures}

\begin{Rblock}
rm_ <- realized_measures(k)
head(rm_[, c("date", "n_5m", "RV", "BV", "Jump")])
\end{Rblock}

This returns realized variance, bipower variation, the jump component,
the two semivariances, realized quarticity, realized skewness and
kurtosis, the Parkinson range, and the activity variables.

Two conventions are worth knowing. Returns are differenced over the full
series rather than within each calendar day, so the interval spanning
midnight is retained and a complete day contributes 288 returns rather
than 287. These markets trade continuously, so excluding it would
discard part of each day. Set \texttt{within\_day = TRUE} to exclude it.
And days with fewer than ninety per cent of the expected intervals have
their measures set to \texttt{NA} rather than being computed from
partial data; change the threshold with \texttt{min\_share}.

\section{Checking what you built}

\begin{Rblock}
max(abs(rm_$RV - (rm_$RSV_pos + rm_$RSV_neg)), na.rm = TRUE)   # ~1e-16
max(abs(rm_$Jump - pmax(rm_$RV - rm_$BV, 0)), na.rm = TRUE)    # ~1e-16
all(rm_$RV > 0, na.rm = TRUE)                                  # TRUE
\end{Rblock}

These identities hold by construction. They are the fastest way to
confirm that a replication has gone right, because three quantities
computed by separate sums over the same returns can only agree to
machine precision if they came from one genuine return series.

% =====================================================================
\chapter{Estimating a HAR model}

\section{Building the features}

\begin{Rblock}
feat <- har_features(rm_$RV, date = rm_$date)
head(feat)
\end{Rblock}

The daily, weekly and monthly components are rolling averages ending at
$t-1$, and \texttt{y} is the value for the next day, so each row pairs a
forecast origin with what it predicts. The logarithm is taken by default,
for a reason explained in Section 5.3.

To add predictors:

\begin{Rblock}
feat <- har_features(rm_$RV, date = rm_$date,
                     extra = rm_[, c("BV", "Jump", "RSV_neg")])
\end{Rblock}

Different averaging windows:

\begin{Rblock}
feat <- har_features(rm_$RV, lags = c(1, 7, 30))
\end{Rblock}

\section{The design matrix}

\begin{Rblock}
des <- har_design(feat)
dim(des$X)

# or a subset of the predictors
des <- har_design(feat, predictors = c("RV_d", "RV_5", "RV_22"))
\end{Rblock}

\section{Why the logarithm}

Realized variance in levels is extremely right skewed. On the
cryptocurrency data used in the paper, skewness runs between 11 and 22
and kurtosis between 194 and 657, and a single day can contribute more
squared error than the entire evaluation sample. A squared error
criterion applied to the level series therefore measures how well a model
happened to do on a handful of days. In logarithms the skewness falls
below 0.7 and the kurtosis below 4.

Forecasts are returned to the level scale where needed, with the
correction $\exp(\hat f + \hat\sigma^2/2)$.

\section{A single fit}

\begin{Rblock}
fit <- tpr_fit(des$X, des$y)
fit
\end{Rblock}

% =====================================================================
\chapter{Rolling forecasts}

\section{The basic call}

\begin{Rblock}
roll <- tpr_rolling(des$X, des$y, window = 365, date = des$date)
summary(roll)
\end{Rblock}

The model is refitted in every window, and the penalty is reselected
every time, so nothing from the future enters any estimate. The output
has one row per forecast origin per estimator, carrying the forecast,
the realization, the penalty, \texttt{q}, the effective degrees of
freedom, the condition number, the in-sample fit and the two fit-gain
quantities.

\section{A path over several penalties}

\begin{Rblock}
path <- tpr_rolling(des$X, des$y, window = 365, rule = "df",
                    target_df = c(3, 2.5, 2, 1.5, 1, 0.5))
plot(path, which = "path")
\end{Rblock}

All six penalties are applied to the \emph{same} window, so the
comparison across them is paired. This matters: if each penalty drew its
own data, differences between estimators would be contaminated by
sampling variation, and you would see least squares appear to vary with
$\lambda$, which is impossible.

\section{Looking at the output}

\begin{Rblock}
plot(roll, which = "forecast")   # realization against forecast
plot(roll, which = "q")          # the scalar over time
plot(roll, which = "lambda")     # the penalty and the effective df
\end{Rblock}

\section{What to expect}

On cryptocurrency data with a data-driven penalty, $q$ typically lies
between 1.01 and 1.06. The two-parameter correction is a one to six per
cent rescaling: real, exactly quantified, but small. It grows as the
penalty tightens, which is why the path in Section 6.2 is the more
informative exercise.

% =====================================================================
\chapter{Evaluating forecasts}

\section{Loss measures}

\begin{Rblock}
w <- data.table::dcast(roll, origin + actual + prev_mean ~ estimator,
                       value.var = "pred")

tpr_metrics(w$actual, w$OLS,        benchmark = w$prev_mean)
tpr_metrics(w$actual, w$`Ridge-I`,  benchmark = w$prev_mean)
tpr_metrics(w$actual, w$`Ridge-II`, benchmark = w$prev_mean)
\end{Rblock}

The benchmark is the prevailing historical mean, which uses only
information available at each forecast origin. A negative out-of-sample
$R^2$ means the model forecasts worse than that simple average.

\section{Diebold and Mariano}

\begin{Rblock}
tpr_dm((w$actual - w$OLS)^2, (w$actual - w$`Ridge-II`)^2)
\end{Rblock}

A positive statistic means the first loss series is larger, so the second
model forecasts better. The long-run variance is estimated with a
Bartlett kernel and a data-driven bandwidth, which matters because
forecast errors from a rolling window are serially correlated.

\section{Model confidence set}

\begin{Rblock}
L <- cbind(OLS = (w$actual - w$OLS)^2,
           RI  = (w$actual - w$`Ridge-I`)^2,
           RII = (w$actual - w$`Ridge-II`)^2)
tpr_mcs(L, alpha = 0.10)
\end{Rblock}

This answers the right question when several models compete: which can
be eliminated, rather than which happens to have the lowest average loss.

\section{Value-at-Risk}

\begin{Rblock}
vf <- w$`Ridge-II`                # variance forecast, level scale
rt <- returns                     # the realized returns for those days

tpr_var_backtest(vf, rt, alpha = 0.01, dist = "normal")
tpr_var_backtest(vf, rt, alpha = 0.01, dist = "student", df = 5)
\end{Rblock}

On cryptocurrency data the Gaussian quantile is rejected at the one per
cent level for almost every model, while a standardised Student-$t$
quantile is not. The variance forecast is not the problem; the assumed
shape of the conditional distribution is.

% =====================================================================
\chapter{The Monte Carlo study}

\section{A short run}

\begin{Rblock}
mc <- tpr_montecarlo(n = c(100, 200), sigma2 = c(0.5, 1.0),
                     reps = 50, seed = 1)
print(mc)
plot(mc, which = "identity")
\end{Rblock}

\section{The full design of the paper}

\begin{Rblock}
mc <- tpr_montecarlo(
  n       = c(100, 200, 400, 500),
  sigma2  = c(0.1, 0.4, 0.5, 0.6, 0.7, 0.8, 1.0, 1.1, 1.2),
  reps    = 1000,
  seed    = 20260911)
\end{Rblock}

This is 36,000 replications and takes some time. Start with a short run
and confirm the verification block before committing to it.

\section{The data-generating process}

$$RV_t = \mu_t \eta_t, \qquad
\eta_t \sim \mathrm{Gamma}(k, 1/k), \quad \mathrm{E}[\eta_t] = 1,$$
with $\mu_t$ the HAR conditional mean. Three properties follow. The
conditional mean of the generated series is exactly the HAR conditional
mean, so the estimated equation is correctly specified rather than an
approximation. Positivity is structural, since $\mu_t \ge \alpha > 0$ and
$\eta_t > 0$, so no truncation or transformation is needed anywhere. And
the regressors are collinear by construction, because the daily component
is one of the terms averaged in the weekly, which is contained in the
monthly. That nesting is a property of the HAR model itself, and it is
the reason shrinkage is worth considering.

\begin{Rblock}
rv <- har_simulate(1000, sigma2 = 0.5)
all(rv > 0)              # TRUE, structurally
mean(rv)                 # approximately 1 with the default parameters
\end{Rblock}

\section{Reading the output}

\begin{Rblock}
print(mc)
\end{Rblock}

The header reports the share of replications in which $q$ exceeded one,
which must be 100 per cent, and the largest Proposition 2 residual, which
must be of order $10^{-16}$. These are identities: a departure indicates
an implementation error, not a property of the data.

\begin{Rblock}
plot(mc, which = "r2")         # out-of-sample fit by sample size
plot(mc, which = "recovery")   # share of sacrificed fit restored
plot(mc, which = "identity")   # empirical against theoretical gain
\end{Rblock}

% =====================================================================
\chapter{A complete worked example}

Everything from download to evaluation, for one asset.

\begin{Rblock}
library(tpridge)
library(data.table)

## 1. download five-minute prices
k <- binance_klines("BTCUSDT", from = "2018-05", to = "2025-12",
                    dir = "~/binance_cache")

## 2. daily realized measures
rm_ <- realized_measures(k)
cat(sprintf("%d days, %s to %s\n", nrow(rm_),
            min(rm_$date), max(rm_$date)))

## 3. check the identities before going further
stopifnot(max(abs(rm_$RV - (rm_$RSV_pos + rm_$RSV_neg)),
              na.rm = TRUE) < 1e-12)

## 4. HAR features with the realized measures as extra predictors
feat <- har_features(rm_$RV, date = rm_$date,
                     extra = rm_[, .(lBV = log(BV), Jump,
                                     lRSVn = log(RSV_neg),
                                     sqRQ = sqrt(RQ))])
des <- har_design(feat)

## 5. rolling forecasts, one year of estimation data
roll <- tpr_rolling(des$X, des$y, window = 365, date = des$date,
                    progress = TRUE)

## 6. how did they do
summary(roll)

## 7. is the difference significant
w <- dcast(roll, origin + actual + prev_mean ~ estimator,
           value.var = "pred")
tpr_dm((w$actual - w$OLS)^2, (w$actual - w$`Ridge-II`)^2)

## 8. did the theory hold in every window
rii <- roll[estimator == "Ridge-II"]
cat(sprintf("q > 1 in %.1f%% of windows\n", 100 * mean(rii$q > 1)))
cat(sprintf("largest Proposition 2 residual %.2e\n",
            max(abs(rii$gap_empirical - rii$gap_theoretical))))

## 9. look at it
plot(roll, which = "forecast")
plot(roll, which = "q")
\end{Rblock}

\section{Doing this for several assets}

\begin{Rblock}
assets <- binance_symbols()
out <- lapply(assets, function(a) {
  k   <- binance_klines(a, "2018-05", "2025-12", dir = "~/binance_cache")
  rm_ <- realized_measures(k)
  des <- har_design(har_features(rm_$RV, date = rm_$date))
  r   <- tpr_rolling(des$X, des$y, window = 365, date = des$date)
  cbind(asset = a, summary(r))
})
rbindlist(out)
\end{Rblock}

% =====================================================================
\chapter{Function reference}

\begin{longtable}{p{0.34\textwidth} p{0.60\textwidth}}
\toprule
\textbf{Estimation} & \\
\midrule
\texttt{tpr\_lm(formula, data)} & Formula interface, for any
regression. Returns an object of class \texttt{tpr\_lm}. \\
\texttt{tpr\_fit(X, y, ...)} & All three estimators from one
decomposition, taking a matrix. \\
\texttt{tpr\_q(X, y, lambda)} & The two-parameter scalar, standalone. \\
\texttt{tpr\_lambda(X, y, rule)} & The penalty, by the data-driven rule
or for a target effective degrees of freedom. \\
\texttt{tpr\_eff\_df(x, lambda)} & Effective degrees of freedom. \\
\texttt{tpr\_window(X, y)} & Decompose a design once, for reuse across
penalties. \\
\addlinespace
\toprule
\textbf{Data} & \\
\midrule
\texttt{binance\_klines(symbol, from, to)} & Download and cache
five-minute archives. \\
\texttt{binance\_symbols()} & The five pairs used in the paper. \\
\texttt{realized\_measures(x)} & Daily measures from intraday prices. \\
\addlinespace
\toprule
\textbf{HAR} & \\
\midrule
\texttt{har\_features(rv, ...)} & Daily, weekly and monthly components
plus the target. \\
\texttt{har\_design(features)} & Design matrix and response, complete
cases. \\
\texttt{har\_simulate(n, sigma2)} & Simulate the process. \\
\addlinespace
\toprule
\textbf{Forecasting} & \\
\midrule
\texttt{tpr\_rolling(X, y, window)} & Rolling one-step forecasts for all
three estimators. \\
\addlinespace
\toprule
\textbf{Evaluation} & \\
\midrule
\texttt{tpr\_metrics(actual, pred)} & Mean squared error, mean absolute
error, out-of-sample $R^2$, quasi-likelihood. \\
\texttt{tpr\_dm(loss\_a, loss\_b)} & Diebold and Mariano with a HAC
long-run variance. \\
\texttt{tpr\_mcs(loss)} & Model confidence set. \\
\texttt{tpr\_var\_backtest(v, r)} & Kupiec and Christoffersen tests. \\
\addlinespace
\toprule
\textbf{Simulation} & \\
\midrule
\texttt{tpr\_montecarlo(...)} & The Monte Carlo design of the paper. \\
\bottomrule
\end{longtable}

\section{Methods}

\texttt{print}, \texttt{summary}, \texttt{coef}, \texttt{fitted},
\texttt{predict} and \texttt{residuals} for \texttt{tpr\_lm};
\texttt{print}, \texttt{coef}, \texttt{fitted}, \texttt{predict} and
\texttt{plot} for \texttt{tpr\_fit}; \texttt{print}, \texttt{summary}
and \texttt{plot} for \texttt{tpr\_rolling}; \texttt{print} and
\texttt{plot} for \texttt{tpr\_mc}.

% =====================================================================
\chapter{Troubleshooting}

\paragraph{``Too few complete observations to estimate the model.''}
The window contains fewer rows than $p+2$ after dropping incomplete
cases. Either the window is too short or a predictor is mostly missing.
Check \texttt{colSums(is.na(X))}.

\paragraph{\texttt{target\_df} must lie in $(0, p]$.}
The effective degrees of freedom cannot exceed the number of predictors.
With three predictors, \texttt{target\_df = 3} recovers least squares and
\texttt{target\_df = 0.5} is heavy shrinkage.

\paragraph{A month fails to download.}
The pair may not have been listed yet. The function reports the month and
continues. \texttt{binance\_klines("XRPUSDT", "2017-01", "2018-12")}
returns data only from May 2018, when the pair was listed.

\paragraph{$q$ is exactly 1.}
This happens when $\lambda = 0$, where the two-parameter estimator
coincides with ridge and with least squares. It is correct, not a
failure.

\paragraph{The Proposition 2 residual is not tiny.}
This should not occur. If it does, the design matrix is probably
singular to working precision. Check
\texttt{tpr\_window(X, y)\$eigenvalues} for values at or below zero.

\paragraph{Out-of-sample $R^2$ exceeds the in-sample $R^2$.}
The two are computed against different benchmarks: the in-sample figure
uses each window's own mean, the out-of-sample figure the prevailing mean
over a longer period. They are not comparable as levels. To compare them,
compute both against the same benchmark.

\paragraph{\texttt{tpr\_mcs()} says the \texttt{MCS} package is required.}
Install it with \texttt{install.packages("MCS")}. It is suggested rather
than required so that the package installs without it.

\paragraph{``lazy-load database is corrupt''.}
Windows will not overwrite files belonging to a package that is
currently loaded, so installing over a version already in use leaves the
database half written. Restart R, then install. If it persists, remove
the directory and any lock first:

\begin{Rblock}
unlink(file.path(.libPaths()[1], c("tpridge", "00LOCK-tpridge")),
       recursive = TRUE, force = TRUE)
## restart R, then install again with
##   devtools::install(".", args = "--no-byte-compile")
\end{Rblock}

\paragraph{``'arg' must be NULL or a character vector''.}
Something in the global environment is masking a package function,
usually an older \texttt{tpr\_fit()} left by sourcing a script.
\texttt{rm(list = ls())}, then reload the package. Turning off
RStudio's restoration of \texttt{.RData} at startup prevents it
recurring.

\paragraph{The installer offers to update dependencies and then fails.}
Decline the update, or restart R before installing so that nothing is
loaded. Installing with

\begin{Rblock}
remotes::install_github("talhaomer472/tpridge", upgrade = "never")
\end{Rblock}

skips the question. The package works with any \texttt{data.table} from
1.14.0 onwards, so there is no reason to force an update.

\paragraph{``Skipping install, the SHA1 has not changed''.}
The installed version already matches the one on GitHub. Nothing is
wrong. Use \texttt{force = TRUE} to reinstall anyway.

\paragraph{Rtools warnings on Windows.}
This package contains no compiled code, so Rtools is not needed. The
warning appears during any source installation and can be ignored.

% =====================================================================
\chapter{References}

\noindent
Christoffersen, P. (1998). Evaluating interval forecasts.
\emph{International Economic Review} 39, 841--862.

\noindent
Corsi, F. (2009). A simple approximate long-memory model of realized
volatility. \emph{Journal of Financial Econometrics} 7, 174--196.

\noindent
Diebold, F. X. and Mariano, R. S. (1995). Comparing predictive accuracy.
\emph{Journal of Business and Economic Statistics} 13, 253--263.

\noindent
Hansen, P. R., Lunde, A. and Nason, J. M. (2011). The model confidence
set. \emph{Econometrica} 79, 453--497.

\noindent
Hoerl, A. E. and Kennard, R. W. (1970). Ridge regression: biased
estimation for nonorthogonal problems. \emph{Technometrics} 12, 55--67.

\noindent
Hoerl, A. E., Kennard, R. W. and Baldwin, K. F. (1975). Ridge
regression: some simulations. \emph{Communications in Statistics} 4,
105--123.

\noindent
Kupiec, P. (1995). Techniques for verifying the accuracy of risk
measurement models. \emph{Journal of Derivatives} 3, 73--84.

\noindent
Longley, J. W. (1967). An appraisal of least squares programs for the
electronic computer from the point of view of the user. \emph{Journal of
the American Statistical Association} 62, 819--841.

\noindent
Lipovetsky, S. and Conklin, W. M. (2005). Ridge regression in
two-parameter solution. \emph{Applied Stochastic Models in Business and
Industry} 21, 525--540.

\end{document}
